%------------------------------------------------------------------------------%
\begin{question}
  Considere os pontos
  \begin{align*}
    A & = ( 1, 1, 0 ) \\
    B & = ( 3, 2, 1 ) \\
    C & = ( 2, 4, 1 ) \\
    D & = ( 1, 1, 3 )
  \end{align*}
  Determine o volume do paralelepípedo que possui $A$ como vértice
  e é determinado pelos vetores
  \[
    \overrightarrow{AB}
    \qquad
    \overrightarrow{AC}
    \qquad
    \overrightarrow{AD}
  \]
\end{question}

%------------------------------------------------------------------------------%
\begin{resolution}{Resolução}
  \[
    \overrightarrow{AB}
    \pause
    = B-A
    \pause
    = \vetortri{2}{1}{1}
  \]

  \[
    \pause
    \overrightarrow{AC}
    \pause
    = C-A
    \pause
    = \vetortri{1}{3}{1}
  \]

  \[
    \pause
    \overrightarrow{AD}
    \pause
    = D-A
    \pause
    = \vetortri{0}{0}{3}
  \]
\end{resolution}

%------------------------------------------------------------------------------%
\begin{resolution}{Resolução}
\begin{columns}[t]
\column{0.5\textwidth}
  \begin{align*}
    \onslide<+->{p}
    \onslide<+->{& = \left(\overrightarrow{AB} \times \overrightarrow{AC}\right) \cdot \overrightarrow{AD}}
    \\
    \onslide<+->{& = \begin{vmatrix}
          2 & 1 & 1 \\
          1 & 3 & 1 \\
          0 & 0 & 3
        \end{vmatrix}}
    \\
    \onslide<+->{& = 3 (-1)^{3+3}
        \begin{vmatrix}
          2 & 1 \\
          1 & 3
        \end{vmatrix}}
    \\
    \onslide<+->{& = 3(6-1)} \\
    \onslide<+->{& = 15}
  \end{align*}
\column{0.5\textwidth}
  \[
    \onslide<+->{V}
    \onslide<+->{= \abs{p}}
    \onslide<+->{= \abs{15}}
    \onslide<+->{= 15}
  \]
\end{columns}
\end{resolution}
%------------------------------------------------------------------------------%
